% Reserve more space for formal result blocks and less for worked material that may flow naturally.
\newcommand{\formalblockneedspace}{10\baselineskip}
\newcommand{\supportingblockneedspace}{7\baselineskip}
\newcommand{\workedblockneedspace}{4\baselineskip}
\newcommand{\proofblockneedspace}{3\baselineskip}
\newcommand{\reserveblockspace}[2]{%
  \AddToHook{env/#1/before}{\Needspace{#2}}%
}
\newcommand{\declareformalblock}[2]{%
  \newtheorem{#1}{#2}[chapter]%
  \reserveblockspace{#1}{\formalblockneedspace}%
}
\newcommand{\declareworkedblock}[2]{%
  \newtheorem{#1}{#2}[chapter]%
  \reserveblockspace{#1}{\workedblockneedspace}%
}
\newcommand{\declaresupportingblock}[2]{%
  \newtheorem{#1}{#2}[chapter]%
  \reserveblockspace{#1}{\supportingblockneedspace}%
}

\theoremstyle{plain}
\declareformalblock{theorem}{Theorem}
\declareformalblock{lemma}{Lemma}
\declareformalblock{proposition}{Proposition}
\declareformalblock{corollary}{Corollary}

\theoremstyle{definition}
\declareformalblock{definition}{Definition}
\declareworkedblock{example}{Example}
\declareworkedblock{exercise}{Exercise}

\theoremstyle{remark}
\declaresupportingblock{remark}{Remark}

\newenvironment{solution}{\begin{proof}[Solution]}{\end{proof}}

\reserveblockspace{proof}{\proofblockneedspace}
